\documentclass[11pt,a4paper]{article}
\usepackage[T1]{fontenc}
\usepackage{lmodern,amsmath,amssymb,amsthm,mathtools}
\usepackage[margin=25mm]{geometry}
\usepackage{microtype,booktabs,enumitem,xcolor,listings,xurl,hyperref,fancyhdr}
\hypersetup{colorlinks=true,linkcolor=blue!45!black,citecolor=blue!45!black,urlcolor=blue!45!black,pdftitle={A Proof of Sun's Divisibility Conjecture for A179524},pdfauthor={Research manuscript prepared with ChatGPT for Keio}}
\setlength{\headheight}{14pt}
\pagestyle{fancy}\fancyhf{}
\fancyhead[L]{Sun's A179524 divisibility conjecture}
\fancyhead[R]{A complete divisibility argument}
\fancyfoot[C]{\thepage}
\newtheorem{theorem}{Theorem}[section]
\newtheorem{lemma}[theorem]{Lemma}
\newtheorem{corollary}[theorem]{Corollary}
\newtheorem{proposition}[theorem]{Proposition}
\theoremstyle{remark}\newtheorem{remark}[theorem]{Remark}
\newcommand{\Z}{\mathbb Z}\newcommand{\Q}{\mathbb Q}
\newcommand{\CT}{\operatorname{CT}}
\newcommand{\Diag}{\operatorname{Diag}}
\lstset{language=Python,basicstyle=\ttfamily\footnotesize,keywordstyle=\color{blue!55!black},commentstyle=\color{gray!80!black},stringstyle=\color{gray!80!black},breaklines=true,columns=fullflexible,keepspaces=true,showstringspaces=false,frame=single,rulecolor=\color{gray!40},aboveskip=10pt,belowskip=10pt}
\title{\textbf{A Proof of Sun's Divisibility\\Conjecture for A179524}\\[5pt]
\large Abel summation, two coefficient states, and integral Catalan coordinates}
\author{Research manuscript prepared with ChatGPT for Keio}
\date{8 November 2026}
\begin{document}\maketitle
\begin{abstract}
Let \(a_n=\sum_k(-4)^k\binom nk^2\binom{n-k}k^2\).
We prove that \(N\) divides \(\sum_{n=0}^{N-1}(20n+17)a_n\) for every
positive integer \(N\), the unchanged divisibility assertion in Sun's
Conjecture 5.8(ii).
We first introduce a diagonal sequence \(Y(N)\). An established
Ramanujan-sum divisibility theorem of Sun, combined with an exact
Abel-summation estimate, proves \(p^e\mid Y(p^e)\) for every odd prime
and every exponent. A digit operator on two copies of the restricted
integral \(p\)-adic power-series ring propagates these bases to all
multiples: the second coefficient state is killed modulo \(p\), giving
a rank-one reduction. The 2-part is termwise. An integral Catalan
coordinate change then converts diagonal divisibility into a
logarithmic divergence, and a constant-term identity returns the result
to the original weighted sum.
All exponents and cofactors are covered, and no finite computation is
an assumption of the proof. The external arithmetic input and the
antecedents of the analytic methods are identified explicitly;
historical priority of the resulting proof is not asserted.
\end{abstract}

\section{The original conjecture and the result}\label{sec:intro}

Set
\begin{equation}\label{eq:original}
 a_n=\sum_{k=0}^{\lfloor n/2\rfloor}
       (-4)^k\binom nk^2\binom{n-k}k^2,
 \qquad S(N)=\sum_{n=0}^{N-1}(20n+17)a_n.
\end{equation}
The initial sequence values are \(1,1,-15,-143,1,12801\).
The assertion \(N\mid S(N)\) is the divisibility part of Sun's
Conjecture 5.8(ii) in \cite[printed p.~33]{Sun2010}, with
\(a_n=W_n(1,-4)\). It is also recorded in OEIS A179524, with its
formulation attributed to 1 July 2010 \cite{OEIS}.

\begin{theorem}\label{thm:main}
For every positive integer \(N\),
\[
                     \boxed{\quad N\mid S(N).\quad}
\]
Consequently, for every prime \(p\equiv1,3,7,9\pmod{20}\),
\[
                   p^e\mid S(p^e)\qquad\text{for all }e\ge1.
\]
\end{theorem}

The proof does not distinguish split and inert primes. Instead it
closes a stronger sufficient condition from the negative-index route
of the supplied manuscript \cite{Supplied} and the integral Catalan
reduction in the preceding note \cite{Previous}. We include that
reduction and the constant-term argument in full, so neither unpublished
note is required as an unproved input.
The only specialized external arithmetic theorem used is Sun's
published Ramanujan partial-sum divisibility result
\cite[Theorem~5, (5)]{Sun2013}.

For clarity, the proof has five steps. First, reverse the diagonal
sum at \(N=p^e\) into a Ramanujan sum with an exact product correction.
Second, use Abel summation to control that correction at every
exponent. Third, construct an integral digit operator on two coefficient
states and propagate the prime-power bases to arbitrary cofactors.
Fourth, handle the 2-part directly. Finally, use integral Catalan
coordinates and formal integration by parts to obtain
Theorem~\ref{thm:main} with the original weight and sequence.

\section{The diagonal sequence}\label{sec:diagonal}

Write \(b_n=\binom{2n}{n}\), and define
\begin{equation}\label{eq:diagonal}
 Y(N)=\sum_{\substack{i,l\ge0\\i+2l=N-1}}
          (-4)^l b_l^2 b_i(10i+3),\qquad N\ge1.
\end{equation}
The auxiliary theorem to be established is
\begin{equation}\label{eq:Ytarget}
                         N\mid Y(N)\qquad(N\ge1).
\end{equation}
For example, \(Y(N)/N\) begins
\(3,13,30,61,506,2962\).
These values only illustrate the definitions; they are not used below.
If
\[
 G(x)=\sum_{l\ge0}(-4)^l b_l^2 x^{2l},
\]
then the binomial series gives
\begin{equation}\label{eq:Ygenerator}
 \sum_{N\ge1}Y(N)x^{N-1}
       =\frac{3+8x}{(1-4x)^{3/2}}G(x).
\end{equation}
Thus \eqref{eq:Ytarget} is also the integrality assertion
\[
 \int_0^x\frac{3+8t}{(1-4t)^{3/2}}G(t)\,dt\in x\Z[[x]].
\]
The implication from \eqref{eq:Ytarget} to the original assertion will
be proved in Sections~\ref{sec:catalan}--\ref{sec:return}; no converse
is required or assumed.



\section{All prime-power diagonal bases by Abel summation}
\label{sec:abel}

Write \(b_j=\binom{2j}{j}\), and set
\[
 A_j=\frac{(20j+3)b_j^2b_{2j}}{(-1024)^j},\qquad
 R_j=\sum_{k=0}^j A_k.
\]
The input to the argument is the following established divisibility theorem
of Sun, not a new result of this paper.

\begin{theorem}[Sun {\cite[Theorem~5, (5)]{Sun2013}}]\label{sun:partial}
For every integer \(j\ge1\),
\[
 4(2j+1)b_j\ \Bigm|\
 \sum_{k=0}^j(20k+3)b_k^2b_{2k}(-1024)^{j-k}
 =(-1024)^jR_j.
\]
\end{theorem}

\begin{lemma}\label{abel:partialbound}
For every odd prime \(p\) and every integer \(j\ge0\),
\[
                  v_p(R_j)\ge v_p\bigl((2j+1)(2j+2)\bigr).
\]
\end{lemma}
\begin{proof}
For \(j\ge1\), Theorem~\ref{sun:partial} and the fact that \(-1024\)
is a \(p\)-adic unit give
\[
 v_p(R_j)\ge v_p(2j+1)+v_p(b_j).
\]
The Catalan number \(b_j/(j+1)\) is an integer, so
\(v_p(b_j)\ge v_p(j+1)=v_p(2j+2)\).
For \(j=0\), the assertion follows directly from \(R_0=3\).
\end{proof}

\begin{theorem}[All odd prime-power diagonal bases]\label{abel:bases}
For every odd prime \(p\) and every integer \(e\ge1\),
\[
                              p^e\mid Y(p^e).
\]
\end{theorem}
\begin{proof}
Put \(N=p^e\) and \(h=(N-1)/2\). For \(0\le j\le h\), define
\[
 r_j=\prod_{i=1}^j
       \left(\frac{1-N/(2i-1)}{1-N/(2i)}\right)^2,
 \qquad r_0=1.
\]
Every factor is a \(p\)-adic unit, since \(1\le2i-1,2i<N\), and hence
\(v_p(N/(2i-1)),v_p(N/(2i))\ge1\).
An exact factorial-ratio identity is
\[
 \left(\frac{b_{h-j}}{b_h}\right)^2
                 =\frac{b_j^2}{256^j}\,r_j.
\]
It follows either by expanding the factorials or by comparing successive
ratios, starting at \(j=0\). In the definition of \(Y(N)\), substitute
\(l=h-j\), so that \(i=N-1-2l=2j\). We obtain the exact identity
\begin{equation}\label{abel:bridge}
 Y(N)=\sum_{j=0}^h(20j+3)(-4)^{h-j}b_{h-j}^2b_{2j}
     =(-4)^h b_h^2\sum_{j=0}^h A_jr_j.
\end{equation}

For \(0\le j<h\), put \(a=2j+1\), \(b=2j+2\). Then
\begin{equation}\label{abel:difference}
 r_{j+1}-r_j
 =-N r_j\frac{2ab-N(a+b)}{a^2(b-N)^2}.
\end{equation}
The positive integers \(a,b\) are consecutive and less than \(N\).
Consequently, if \(t=v_p(ab)\), then \(t<e\),
\(v_p(b-N)=v_p(b)\), and
\[
 v_p\bigl(2ab-N(a+b)\bigr)=t.
\]
Here the first summand has valuation \(t\), because \(p\) is odd,
whereas the second has valuation at least \(e>t\).
As \(r_j\) is a unit, \eqref{abel:difference} gives
\begin{equation}\label{abel:valuation}
 v_p(r_{j+1}-r_j)=e-v_p\bigl((2j+1)(2j+2)\bigr).
\end{equation}

Abel summation yields
\[
 \sum_{j=0}^h A_jr_j
       =R_hr_h-\sum_{j=0}^{h-1}R_j(r_{j+1}-r_j).
\]
Every summand in the last sum belongs to \(p^e\Z_p\), by
Lemma~\ref{abel:partialbound} and \eqref{abel:valuation}.
Theorem~\ref{sun:partial} at \(h\ge1\) also gives
\(R_h\in p^e\Z_p\), because \(2h+1=N\).
Thus the entire right-hand side is in \(p^e\Z_p\).
The prefactor in \eqref{abel:bridge} is an integer, and \(Y(N)\) is an
integer, proving the asserted divisibility.
\end{proof}

\section{Exact integral kernels for central binomial coefficients}
\label{sec:central-kernels}

Throughout this section, \(p\) is an arbitrary odd prime.
Let \(\Z_p\langle X_1,\ldots,X_d\rangle\) denote the ring of restricted
power series: its coefficients lie in \(\Z_p\) and tend to zero as total
degree tends to infinity. It is complete for the Gauss norm. Integral
affine substitutions preserve this ring, and its elements may be evaluated
at every point of \(\Z_p^d\).

\begin{lemma}[Normalized unit-factorial interpolation]\label{kernel:units}
Put
\[
 E_p(M)=\prod_{\substack{1\le i\le M\\p\nmid i}}i,
 \qquad W=(p-1)!.
\]
For each \(0\le r<p\), there is a unit
\(\mathcal H_r(X)\in1+p\Z_p\langle X\rangle\), with
\(\mathcal H_r(0)=1\), such that
\[
 E_p(pq+r)=W^q r!\,\mathcal H_r(q)\qquad(q\ge0).
\]
\end{lemma}
\begin{proof}
Set
\[
 B(X)=\prod_{i=1}^{p-1}(1+pX/i),\qquad
 \log B(X)=\sum_{j\ge1}\ell_jX^j,\qquad
 \ell_j=\frac{(-1)^{j-1}p^j}{j}\sum_{i=1}^{p-1}i^{-j}.
\]
The polynomial continuation of \(\sum_{h=0}^{q-1}h^j\) is
\[
 S_j(X)=\sum_{a=0}^j
 \left\{\begin{matrix}j\\a\end{matrix}\right\}
 \frac{(X)_{\underline{a+1}}}{a+1},
\]
where braces denote Stirling numbers of the second kind and the underlined
subscript denotes a falling factorial. Thus
\[
 S_j(X)\in p^{-\lfloor\log_p(j+1)\rfloor}\Z_p[X].
\]
For every odd \(p\) and \(j\ge1\),
\[
 j-v_p(j)-\lfloor\log_p(j+1)\rfloor\ge1,
\]
and this lower bound tends to infinity with \(j\).
For completeness, check \(j=1,2,3\) directly; for \(j\ge4\) use
\(v_p(j)\le\lfloor\log_3(j+1)\rfloor\) and
\(2\lfloor\log_3(j+1)\rfloor\le j-1\).
The latter is direct for \(4\le j\le7\); for
\(t=\lfloor\log_3(j+1)\rfloor\ge2\), use
\(j\ge3^t-1\) and \(2t\le3^t-2\).
It follows that
\[
 A(X)=\sum_{j\ge1}\ell_jS_j(X)\in p\Z_p\langle X\rangle.
\]
Because \(p\) is odd, \(\exp A(X)\) converges in
\(1+p\Z_p\langle X\rangle\): the \(k\)-th exponential term has
valuation at least \(k-v_p(k!)\), which tends to infinity.
Define
\[
 \mathcal H_r(X)=\exp A(X)\prod_{i=1}^r(1+pX/i).
\]
All \(S_j(0)\) vanish, so \(\mathcal H_r(0)=1\).
For a nonnegative integer \(q\), convergence and the finite sum over
\(0\le h<q\) give
\[
 A(q)=\sum_{h=0}^{q-1}\log B(h),\qquad
 \exp A(q)=\prod_{h=0}^{q-1}B(h).
\]
Splitting the product \(E_p(pq+r)\) into its \(q\) complete blocks and
its final \(r\) factors proves the asserted identity.
Only the normalized factor is interpolated; no interpolation of \(W^q\)
is required.
\end{proof}

\begin{lemma}[Central-binomial digit kernel]\label{kernel:central}
For \(0\le r<p\), put
\[
 \delta_r=\left\lfloor\frac{2r}{p}\right\rfloor\in\{0,1\},
 \qquad t_r=2r-p\delta_r.
\]
Define
\[
 \mathcal B_r(X)=
 p^{\delta_r}(2X+1)^{\delta_r}W^{\delta_r}
 \frac{t_r!}{r!^2}
 \frac{\mathcal H_{t_r}(2X+\delta_r)}{\mathcal H_r(X)^2}.
\]
Then \(\mathcal B_r\in p^{\delta_r}\Z_p\langle X\rangle\), and
\[
 b_{pq+r}=b_q\mathcal B_r(q)\qquad(q\ge0).
\]
Moreover,
\[
 \mathcal B_r(0)=b_r,\qquad
 \mathcal B_r(X)\equiv b_r\pmod p.
\]
\end{lemma}
\begin{proof}
All the factorial denominators are units, and the \(\mathcal H\)'s are
units in the restricted ring. Use the exact identity
\[
 M!=p^{\lfloor M/p\rfloor}\lfloor M/p\rfloor!E_p(M)
\]
in \(b_{pq+r}\), and apply Lemma~\ref{kernel:units}. The factorial part
remaining after the unit factorials have been removed is
\[
 p^{\delta_r}\frac{(2q+\delta_r)!}{q!^2}
       =p^{\delta_r}b_q(2q+1)^{\delta_r}.
\]
The powers of \(W\) cancel to \(W^{\delta_r}\), proving the exact formula.
Taking \(q=0\) gives the constant value.
If \(\delta_r=1\), both sides of the final congruence vanish modulo \(p\).
If \(\delta_r=0\), reduce the unit factors modulo \(p\) and use
\(t_r=2r\).
\end{proof}

Let
\[
 w(i,l)=(-4)^l b_i b_l^2,\qquad
 \mathcal A_p=\Z_p\langle X,Y\rangle.
\]
Since \((-4)^{p-1}\in1+p\Z_p\), the restricted series
\[
 \mathcal C(Y)=\exp\bigl(Y\log((-4)^{p-1})\bigr)
\]
is a unit in \(1+p\Z_p\langle Y\rangle\), and
\(\mathcal C(j)=(-4)^{(p-1)j}\) for every integer \(j\ge0\).
Define
\begin{equation}\label{kernel:double}
 K_{r,s}(X,Y)=(-4)^s\mathcal C(Y)
                      \mathcal B_r(X)\mathcal B_s(Y)^2.
\end{equation}
Lemma~\ref{kernel:central} gives, for all \(q,j\ge0\),
\begin{equation}\label{kernel:exact}
 w(pq+r,pj+s)=w(q,j)K_{r,s}(q,j).
\end{equation}
In addition,
\begin{equation}\label{kernel:properties}
 \begin{gathered}
 K_{r,s}\in p^{\delta_r+2\delta_s}\mathcal A_p,\qquad
 K_{r,s}(0,0)=w(r,s),\\
 K_{r,s}(X,Y)\equiv w(r,s)\pmod p.
 \end{gathered}
\end{equation}
These formulas include every carry in the two central binomial coefficients.

\section{A two-state coefficient operator}
\label{sec:two-state}

The two neighboring coefficient indices below form a closed system under
base-\(p\) digit decomposition. This is the additional propagation mechanism
needed to use Theorem~\ref{abel:bases}.

For \(F=(F_0,F_1)\in\mathcal M_p:=\mathcal A_p^2\) and \(N\ge1\), set
\begin{equation}\label{state:functional}
 \mathcal J_N(F)=
 \sum_{\substack{i,l\ge0\\i+2l=N-1}}w(i,l)F_0(i,l)
 +\sum_{\substack{i,l\ge0\\i+2l=N-2}}w(i,l)F_1(i,l).
\end{equation}
A sum with negative right-hand side is empty. In particular,
\(\mathcal J_1(F)=F_0(0,0)\).

\begin{proposition}[Exact digit decomposition]\label{state:operator}
For \(d=0,1\), define
\[
 (\mathcal L_pF)_d(X,Y)=
 \sum_{a=0}^1
 \sum_{\substack{0\le r,s<p\\r+2s=p-1-a+pd}}
 K_{r,s}(X,Y)F_a(pX+r,pY+s).
\]
Then \(\mathcal L_p\) is a \(\Z_p\)-linear operator on \(\mathcal M_p\),
and for every \(N\ge1\),
\begin{equation}\label{state:block}
               \mathcal J_{pN}(F)=\mathcal J_N(\mathcal L_pF).
\end{equation}
For every \(F\in\mathcal M_p\), there is \(R_F\in\mathcal M_p\) such that
\begin{equation}\label{state:rank-one}
 \mathcal L_pF=(\mathcal J_p(F),0)+pR_F,
 \qquad (R_F)_0(0,0)=0.
\end{equation}
\end{proposition}
\begin{proof}
Write \(i=pq+r\), \(l=pj+s\) in the summand with input index \(a\in\{0,1\}\).
The constraint \(i+2l=pN-1-a\) becomes
\[
 r+2s=p-1-a+pd,\qquad q+2j=N-1-d.
\]
The only possible values of \(d\) are \(0,1\): the low sum lies between
\(0\) and \(3p-3\), whereas \(d=-1\) would make it negative and
\(d=2\) would make it at least \(3p-2\).
Equation~\eqref{kernel:exact} now proves \eqref{state:block}, including
all empty boundary ranges. The asserted operator preserves
\(\mathcal M_p\) because it is a finite sum of integral restricted-series
products and integral affine substitutions.

For \(d=1\), the low sum is \(2p-1\) or \(2p-2\).
If \(\delta_r=\delta_s=0\), then
\[
 r+2s\le\frac{3(p-1)}2<2p-2,
\]
a contradiction. Hence \(\delta_r+2\delta_s\ge1\), and
\eqref{kernel:properties} shows that
\((\mathcal L_pF)_1\in p\mathcal A_p\).

For \(d=0\), use \eqref{kernel:properties} and
\[
 F_a(pX+r,pY+s)\equiv F_a(r,s)\pmod p.
\]
The latter congruence follows for polynomials by expansion, and for
restricted series by Gauss-norm continuity. We obtain
\[
 (\mathcal L_pF)_0\equiv
 \sum_{a=0}^1\sum_{\substack{0\le r,s<p\\r+2s=p-1-a}}
          w(r,s)F_a(r,s)=\mathcal J_p(F)\pmod p.
\]
The constant-term equality is exact, because
\(K_{r,s}(0,0)=w(r,s)\):
\[
                 (\mathcal L_pF)_0(0,0)=\mathcal J_p(F).
\]
Subtracting the constant vector and dividing both components by \(p\)
gives \eqref{state:rank-one}, including the specified zero constant term.
\end{proof}

\begin{theorem}[Propagation from a chain of prime-power bases]\label{state:propagation}
Let \(p\) be an odd prime, let \(F\in\mathcal M_p\), and let \(e\ge1\).
Suppose
\[
 \mathcal J_{p^j}(F)\in p^j\Z_p\qquad(1\le j<e).
\]
Then, for every \(m\ge1\),
\[
 \mathcal J_{p^em}(F)\equiv
 \mathcal J_{p^e}(F)\,\mathcal J_m(1,0)\pmod{p^e}.
\]
If the base divisibility also holds at \(j=e\), then
\(\mathcal J_{p^em}(F)\in p^e\Z_p\) for every \(m\ge1\).
\end{theorem}
\begin{proof}
Iterating \eqref{state:block}, and using \(\mathcal J_1(F)=F_0(0,0)\), gives
\[
                 (\mathcal L_p^jF)_0(0,0)=\mathcal J_{p^j}(F).
\]
Suppose the base divisibilities have been used through \(j-1\), so that
\(\mathcal L_p^{j-1}F=p^{j-1}G\) with \(G\in\mathcal M_p\).
For \(j=1\), this means simply \(G=F\).
By \eqref{state:rank-one},
\[
 \mathcal L_p^jF
 =p^{j-1}(\mathcal J_p(G),0)+p^jR_G
 =(\mathcal J_{p^j}(F),0)+p^jR_G.
\]
The second equality follows by evaluating the first component at \((0,0)\),
since \((R_G)_0(0,0)=0\).
When the next base is divisible by \(p^j\), the entire iterate belongs
to \(p^j\mathcal M_p\). Induction through \(e-1\), followed by one more
application of the formula, therefore gives
\[
 \mathcal L_p^eF=(\mathcal J_{p^e}(F),0)+p^eR
 \qquad(R\in\mathcal M_p).
\]
Apply \(\mathcal J_m\), whose values on integral pairs are integral, to
obtain the asserted congruence and its divisibility consequence.
The hypothesis is empty when \(e=1\), and the same proof applies.
\end{proof}

\begin{corollary}[Every odd primary part of the diagonal condition]\label{state:all-odd}
For every odd prime \(p\) and all integers \(e,m\ge1\),
\[
                              p^e\mid Y(p^em).
\]
\end{corollary}
\begin{proof}
Take \(F=(10X+3,0)\). Then \(\mathcal J_N(F)=Y(N)\).
All its prime-power bases are supplied by Theorem~\ref{abel:bases}.
Apply Theorem~\ref{state:propagation}.
\end{proof}

\section{The full diagonal integrality theorem}
\label{sec:diagonal-complete}

\begin{lemma}[Termwise divisibility at the prime two]\label{diag:two}
For every \(N\ge1\), \(2^{v_2(N)}\mid Y(N)\).
\end{lemma}
\begin{proof}
Write \(s_2(a)\) for the binary digit sum of a nonnegative integer \(a\).
Legendre's factorial formula gives \(v_2(b_a)=s_2(a)\).
For a summand indexed by \(i+2l=N-1\), the factor \(10i+3\) is odd, so
\[
 v_2\bigl((-4)^l b_l^2b_i(10i+3)\bigr)
                         =2l+2s_2(l)+s_2(i).
\]
Digit-sum subadditivity yields
\[
 s_2(i)+s_2(l)=s_2(i)+s_2(2l)
                     \ge s_2(N-1)\ge v_2(N).
\]
The final inequality holds because the last \(v_2(N)\) binary digits
of \(N-1\) are ones. Thus each individual summand is divisible by
\(2^{v_2(N)}\), and so is their sum. The case \(N=1\) is included,
with exponent zero.
\end{proof}

\begin{theorem}[Integral diagonal primitive]\label{diag:complete}
For every positive integer \(N\),
\[
                                \boxed{N\mid Y(N)}.
\]
Equivalently, with
\(G(x)=\sum_{l\ge0}(-4)^l b_l^2x^{2l}\),
\[
 \int_0^x\frac{3+8t}{(1-4t)^{3/2}}G(t)\,dt\in\Z[[x]].
\]
\end{theorem}
\begin{proof}
For every odd prime divisor \(p\) of \(N\), write
\(N=p^em\), where \(e=v_p(N)\ge1\).
Corollary~\ref{state:all-odd} gives \(p^e\mid Y(N)\).
Lemma~\ref{diag:two} supplies the entire \(2\)-part.
These prime-power divisors are pairwise coprime, so their product \(N\)
divides \(Y(N)\). For \(N=1\), the assertion is immediate.
The integrand in the displayed formal antiderivative is
\(\sum_{N\ge1}Y(N)x^{N-1}\), proving the equivalent formulation.
\end{proof}


\section{An integral Catalan criterion}\label{sec:catalan}

We now give the algebraic bridge from the proved diagonal divisibility
to the original sum. Everything in this section takes place over
\(\Z\); there is no convergence condition beyond formal power series.
Let
\[
 A=\Z[[u,v]],\qquad
 \theta_u=u\frac{\partial}{\partial u},\quad
 \theta_v=v\frac{\partial}{\partial v},\qquad
 \mathcal D_A=\theta_u A+\theta_v A.
\]
A coefficientwise application of B\'ezout's identity shows that
\begin{equation}\label{eq:bezout}
 F\in\mathcal D_A
 \quad\Longleftrightarrow\quad
 \begin{cases}
 [u^0v^0]F=0,\\
 \gcd(r,s)\mid[u^rv^s]F\quad((r,s)\ne(0,0)).
 \end{cases}
\end{equation}
Indeed, the coefficient of \(u^rv^s\) in \(\theta_u H+\theta_v K\)
is \(rH_{r,s}+sK_{r,s}\), and these integer combinations form precisely
\(\gcd(r,s)\Z\).

\begin{theorem}[Diagonal criterion]\label{thm:catalan}
Let \(R(x)\in\Z[[x]]\), \(\alpha,\beta\in\Z\), and
\[
 \rho=\frac{uv}{1+u+v},\qquad
 F(u,v)=(\alpha\rho+\beta\rho^2)R(uv).
\]
Then
\[
 F\in\mathcal D_A
 \quad\Longleftrightarrow\quad
 n\mid[u^nv^n]F\qquad\text{for every }n\ge1.
\]
\end{theorem}
\begin{proof}
The forward direction follows from \eqref{eq:bezout}.
For the converse let \(C(x)\in1+x\Z[[x]]\) be the Catalan series,
characterized by \(C=1+xC^2\), and set
\[
 U=uC(uv),\qquad V=vC(uv),\qquad y=UV.
\]
This is an integral formal automorphism. Its inverse and the induced
radial substitution are
\begin{equation}\label{eq:catalaninverse}
 u=\frac{U}{1+y},\qquad v=\frac{V}{1+y},\qquad
 x=uv=\frac{y}{(1+y)^2},\qquad C(x)=1+y.
\end{equation}
The inverse logarithmic differential matrix is
\[
 \begin{pmatrix}d\log u\\d\log v\end{pmatrix}
 =\frac1{1+y}
  \begin{pmatrix}1&-y\\-y&1\end{pmatrix}
  \begin{pmatrix}d\log U\\d\log V\end{pmatrix}.
\]
Its determinant
\[
                         J(y)=\frac{1-y}{1+y}
\]
is an integral unit, and the matrix and its inverse have integral
power-series entries. In the free module of logarithmic differentials,
\[
 F\,d\log u\wedge d\log v
       =d(H\,d\log v-K\,d\log u)
 \quad\Longleftrightarrow\quad F=\theta_uH+\theta_vK.
\]
Thus pullback preserves integral exactness in both directions. In
particular \(F\in\mathcal D_A\) if and only if
\[
 \widetilde F(U,V)
       =J(y)F\left(\frac{U}{1+y},\frac{V}{1+y}\right)
\]
is an integral logarithmic divergence in \(U,V\).

The useful factorization is
\[
 1+u+v=\frac{(1+U)(1+V)}{1+y},\qquad
 \rho=\frac{y}{(1+y)(1+U)(1+V)}.
\]
It gives
\begin{equation}\label{eq:separated}
 \widetilde F
   =\frac{A_1(y)}{(1+U)(1+V)}
    +\frac{A_2(y)}{(1+U)^2(1+V)^2},
\end{equation}
where
\[
 A_1(y)=\frac{\alpha y(1-y)}{(1+y)^2}R\left(\frac{y}{(1+y)^2}\right),
 \qquad
 A_2(y)=\frac{\beta y^2(1-y)}{(1+y)^3}R\left(\frac{y}{(1+y)^2}\right).
\]
These are integral; \(A_1\) starts in degree one and \(A_2\) in degree
two. Write \(A_j(y)=\sum_k A_{j,k}y^k\). For \(r,s\ge1\), set
\(q=\min(r,s)\) and \(d=\gcd(r,s)\). Expanding the denominators in
\eqref{eq:separated},
\[
 [U^rV^s]\widetilde F
   =(-1)^{r+s}\sum_{k=0}^{q}
       \bigl(A_{1,k}+A_{2,k}(r-k+1)(s-k+1)\bigr).
\]
Since \(r,s,q\) are all multiples of \(d\),
\begin{equation}\label{eq:gcdreduction}
 [U^rV^s]\widetilde F
       \equiv(-1)^{r+s}[U^qV^q]\widetilde F\pmod d.
\end{equation}

We must also preserve the diagonal divisibility under the change of
variables. Let \(\Diag F(x)=\sum_{n\ge0}[u^nv^n]F\,x^n\).
The radial substitution \eqref{eq:catalaninverse} preserves the
difference of the two exponents, hence
\[
 \Diag\widetilde F(y)=J(y)\,\Diag F\left(\frac{y}{(1+y)^2}\right).
\]
The assumed diagonal divisibility gives an integral series
\(H_0(x)\in x\Z[[x]]\) such that \(\Diag F(x)=xH_0'(x)\).
For \(x(y)=y/(1+y)^2\), the identity
\(yx'(y)/x(y)=J(y)\) therefore gives
\[
              \Diag\widetilde F(y)
                       =y\frac{d}{dy}H_0(x(y)).
\]
Its coefficient of \(y^q\) is divisible by \(q\).
Equation~\eqref{eq:gcdreduction} now proves every required gcd
condition. The constant and axis coefficients vanish because of the
initial degrees of \(A_1,A_2\). By \eqref{eq:bezout}, \(\widetilde F\)
is an integral logarithmic divergence. Pulling back by the inverse
automorphism proves the same for \(F\).
\end{proof}

\begin{corollary}\label{cor:Psi}
Let \(G(x)\) be as in Section~\ref{sec:diagonal}, and put
\[
                       \Psi=(3\rho+20\rho^2)G(uv).
\]
There exist \(H,K\in\Z[[u,v]]\) with
\[
                            \Psi=\theta_uH+\theta_vK.
\]
\end{corollary}
\begin{proof}
For \(i\ge0\),
\[
 [u^{i+1}v^{i+1}]\rho=\binom{2i}{i}.
\]
For \(i\ge1\),
\[
 [u^{i+1}v^{i+1}]\rho^2
       =(2i-1)\binom{2i-2}{i-1},\qquad
 20(2i-1)\binom{2i-2}{i-1}=10i\binom{2i}{i}.
\]
For \(i=0\) the \(\rho^2\) coefficient is zero. Consequently
\([u^Nv^N]\Psi=Y(N)\) for every \(N\ge1\).
The already proved divisibility \eqref{eq:Ytarget} and
Theorem~\ref{thm:catalan} give the assertion.
\end{proof}

\section{Returning to the original weighted sum}\label{sec:return}

Put
\[
 \Lambda=\frac{(1+u)(1+v)}{uv},\qquad
 B=\bigcup_{M\ge0}(uv)^{-M}\Z[[u,v]].
\]
Constant terms of the expressions below are well defined in \(B\):
all exponents have a fixed lower bound in each variable. The formal
series in \(A\) have integer coefficients, so these constant terms
are integers. For every \(n\ge0\),
\begin{equation}\label{eq:ct-original}
 \CT\bigl(\Lambda^nG(uv)\bigr)
    =\sum_{l\ge0}(-4)^l\binom{2l}{l}^2\binom{n}{2l}^2
    =a_n.
\end{equation}
We used \(\binom{n}{2l}\binom{2l}{l}
=\binom nl\binom{n-l}{l}\), including the zero convention outside
the admissible range.

\begin{lemma}[Exact constant-term identity]\label{lem:negative}
For every \(N\ge1\),
\begin{equation}\label{eq:negative}
 S(N)=-\CT(\Lambda^N\Psi)
                   +20N\,\CT\bigl(\Lambda^N\rho G(uv)\bigr).
\end{equation}
\end{lemma}
\begin{proof}
An exact polynomial identity is
\[
 (T-1)^2\sum_{n=0}^{N-1}(20n+17)T^n
    =(3T+17)(1-T^N)+20NT^N(T-1).
\]
Substitute \(T=\Lambda\) and use
\(\rho=(\Lambda-1)^{-1}=uv/(1+u+v)\). Since
\((3\Lambda+17)\rho^2=3\rho+20\rho^2\),
\[
 \sum_{n=0}^{N-1}(20n+17)\Lambda^n
    =(3\rho+20\rho^2)(1-\Lambda^N)+20N\Lambda^N\rho.
\]
Multiply by \(G(uv)\), take the constant term, and apply
\eqref{eq:ct-original}. The remaining term \(\CT\Psi\) vanishes
because \(\Psi\in uvA\).
\end{proof}

\begin{proof}[Proof of Theorem~\ref{thm:main}]
By Corollary~\ref{cor:Psi}, write
\(\Psi=\theta_uH+\theta_vK\) with \(H,K\in A\).
The constant term of a logarithmic derivative in \(B\) is zero.
Since
\[
 \frac{\theta_u\Lambda}{\Lambda}=-\frac1{1+u},\qquad
 \frac{\theta_v\Lambda}{\Lambda}=-\frac1{1+v},
\]
formal integration by parts gives
\[
 \CT(\Lambda^N\Psi)
       =N\,\CT\left[\Lambda^N
               \left(\frac{H}{1+u}+\frac{K}{1+v}\right)\right]
       \in N\Z.
\]
The second term of \eqref{eq:negative} also lies in \(N\Z\).
Thus \(N\mid S(N)\) for every positive integer \(N\), as claimed.
\end{proof}

\begin{corollary}\label{cor:allcofactors}
For every prime \(p\) and all \(e,m\ge1\),
\[
                            p^e\mid S(p^em).
\]
In particular the entire split-prime-power family
\(p\equiv1,3,7,9\pmod{20}\) is included.
\end{corollary}
\begin{proof}
Apply Theorem~\ref{thm:main} with \(N=p^em\).
\end{proof}


\section{Scope, provenance, and verification}\label{sec:scope}

The result proved here is exactly the all-integer divisibility assertion
in Conjecture 5.8(ii) of \cite{Sun2010}. The same part of that conjecture
also predicts the more precise congruence
\[
 S(p)\equiv p\left(10\left(\frac{-1}{p}\right)+7\right)\pmod{p^2}
 \qquad(p\text{ odd prime}).
\]
This congruence is not a conclusion of the present paper. Nor do we
claim to prove the separate unweighted CM evaluations in part (i).
Theorem~\ref{thm:main} covers the exact original sequence, original
weight, all positive integers, and consequently every requested split
prime power.

The external arithmetic input is Theorem~\ref{sun:partial}, a published
theorem of Sun. The additional arguments connecting it to the target are
the exact reflected product in \eqref{abel:bridge}, the valuation in
\eqref{abel:valuation}, the closed pair of coefficient states in
Proposition~\ref{state:operator}, and the integral diagonal criterion
of Theorem~\ref{thm:catalan}. In particular, a formal identity alone
does not prove the target: the Abel corrections must lie in \(p^e\Z_p\),
and the digit operator must preserve an integral coefficient module
at every iteration. Those two assertions are proved in detail above.

The use of normalized factorial units, digit decompositions, and
\(p\)-adic coefficient estimates has substantial antecedents. For
related arithmetic methods for Ap\'ery-like sequences, see Delaygue
\cite{Delaygue}. The present paper does not assert that such general
principles originate here. The supplied research draft \cite{Supplied}
provided the negative-index direction, and the preceding research note
\cite{Previous} supplied the integral Catalan criterion used as a
starting point. All unpublished ingredients needed for the conclusion
have been proved again in this paper.

The bibliography and the accompanying source review are bounded.
A search using the exact sequence identifier, the original weight,
equivalent binomial notation, and the auxiliary convolution did not
identify an earlier complete proof of this exact assertion or of the
full chain presented here.
Neither the appearance of a conjecture in OEIS nor a failure to locate
an identical proof establishes historical priority. Accordingly, this
manuscript makes no claim that the theorem or its essential proof has
never appeared elsewhere. The work is a newly prepared proof manuscript;
its internal checking is not external peer review or a proof-assistant
verification. Mathematical completeness of the argument and a certified
priority claim are separate questions.

The finite tests in the appendix are transcription checks. They use
exact integers and rational numbers, and check the original sums, the
reflected-product identity, the Abel correction valuations, the diagonal
coefficients, and the final constant-term identity. In the accompanying
run, all assertions passed: \(1\le N\le400\) for the original and
diagonal sums, eleven prime-power bases for the Abel identities,
\(1\le r,s\le40\) for the coefficient gcd conditions, and
\(1\le N\le30\) for the constant-term identity.
No upper bound from a computation is used anywhere in the proof.

\appendix
\section{Exact reproducibility checks}\label{app:checks}

The following complete Python program uses only the standard library.
It is included in this single source file for reproducibility.
Its finite assertions supplement, but do not replace, the quantified
proofs in the main text.
\begin{lstlisting}


"""Finite transcription checks; no assertion below replaces a proof."""
from fractions import Fraction as Q
from functools import lru_cache
from math import comb, gcd

@lru_cache(None)
def b(j):
    return comb(2*j, j)

def a(n):
    return sum((-4)**k * comb(n, k)**2 * comb(n-k, k)**2
               for k in range(n//2 + 1))


def Y(n):
    return sum((-4)**l * b(l)**2 * b(n-1-2*l)
               * (10*(n-1-2*l)+3)
               for l in range((n-1)//2 + 1))


def vp(x, p):
    x = Q(x)
    if x == 0:
        return float('inf')
    z, v = abs(x.numerator), 0
    while z % p == 0:
        z //= p
        v += 1
    z = x.denominator
    while z % p == 0:
        z //= p
        v -= 1
    return v


def rho_coefficient(r, s, power):
    if min(r, s) < power:
        return 0
    d = r+s-2*power
    return (-1)**d * comb(power+d-1, d) * comb(d, r-power)


@lru_cache(None)
def coeff(r, s, kind):
    return sum((-4)**l * b(l)**2 * (
        (3 if kind == 'psi' else 1)
        * rho_coefficient(r-2*l, s-2*l, 1)
        + (20 if kind == 'psi' else 0)
        * rho_coefficient(r-2*l, s-2*l, 2))
        for l in range((min(r, s)-1)//2 + 1))


def ct(n, kind):
    return sum(comb(n, r)*comb(n, s)*coeff(r, s, kind)
               for r in range(n+1) for s in range(n+1))


partial = [0]
for n in range(1, 401):
    partial.append(partial[-1] + (20*(n-1)+17)*a(n-1))
    assert partial[n] % n == 0
    assert Y(n) % n == 0
assert [Y(n)//n for n in range(1, 7)] == [3, 13, 30, 61, 506, 2962]

bases = [(3, 1), (3, 2), (3, 3), (3, 4), (5, 1), (5, 2),
         (5, 3), (7, 1), (7, 2), (11, 2), (13, 2)]
for p, e in bases:
    n, r, R, total = p**e, Q(1), Q(0), Q(0)
    h = (n-1)//2
    for j in range(h+1):
        A = Q((20*j+3)*b(j)**2*b(2*j), (-1024)**j)
        R += A
        total += A*r
        assert Q(b(h-j), b(h))**2 == Q(b(j)**2, 256**j)*r
        assert vp(R, p) >= vp((2*j+1)*(2*j+2), p)
        if j < h:
            aa, bb = 2*j+1, 2*j+2
            next_r = r*(Q(1)-Q(n, aa))**2/(Q(1)-Q(n, bb))**2
            delta = next_r-r
            assert delta == -n*r*Q(2*aa*bb-n*(aa+bb), aa**2*(bb-n)**2)
            assert vp(delta, p) == e-vp(aa*bb, p)
            assert vp(R*delta, p) >= e
            r = next_r
    assert Y(n) == (-4)**h*b(h)**2*total
    assert vp(R, p) >= e

for r in range(1, 41):
    for s in range(1, 41):
        assert coeff(r, s, 'psi') % gcd(r, s) == 0
for n in range(1, 31):
    assert coeff(n, n, 'psi') == Y(n)
    assert partial[n] == -ct(n, 'psi') + 20*n*ct(n, 'rho')
print('PASS: original and diagonal divisibility for 1 <= n <= 400')
print('PASS: exact Abel identities and valuations at 11 prime powers')
print('PASS: 1600 coefficient gcd tests and 30 constant-term identities')


\end{lstlisting}

\begin{thebibliography}{9}\raggedright

\bibitem{Sun2010}
Zhi-Wei Sun,
\emph{On Ap\'ery numbers and generalized central trinomial coefficients},
arXiv:1006.2776v11, 19 August 2010, Conjecture 5.8(ii), printed p.~33.
\url{https://arxiv.org/abs/1006.2776v11}.
The version identifier is specified because a later version was
withdrawn when the material was reorganized.

\bibitem{Sun2013}
Zhi-Wei Sun,
\emph{Products and sums divisible by central binomial coefficients},
Electronic Journal of Combinatorics \textbf{20}(1) (2013), Paper P9,
14 pp.; Theorem~5, equation~(5).
\url{https://www.combinatorics.org/ojs/index.php/eljc/article/download/v20i1p9/pdf/}.

\bibitem{OEIS}
OEIS Foundation Inc.,
\emph{The On-Line Encyclopedia of Integer Sequences}, entry A179524,
contributed by Zhi-Wei Sun, 17 July 2010; accessed 8 November 2026.
\url{https://oeis.org/A179524}.

\bibitem{Delaygue}
\'Eric Delaygue,
\emph{Arithmetic properties of Ap\'ery-like numbers},
Compositio Mathematica \textbf{154} (2018), 249--274.
\url{https://doi.org/10.1112/S0010437X17007552}.
Preprint: \url{https://arxiv.org/abs/1310.4131}.

\bibitem{Supplied}
QEEDDD / Verifiable Mathematical Reasoning Corpus project,
\emph{Sun's divisibility conjecture for
\(\sum_{k<n}(20k+17)W_k(1,-4)\): squarefree moduli, the second layer,
and a reduction via negative-index continuation},
user-supplied unpublished research manuscript, 7 November 2026,
13 pp.

\bibitem{Previous}
Research manuscript prepared with ChatGPT for Keio,
\emph{Prime-Power Divisibility in Sun's A179524 Conjecture},
unpublished research note, 8 November 2026, 23 pp.
The present paper reproduces all required arguments and supplies
the remaining prime-power bases and their propagation.

\end{thebibliography}
\end{document}
